There is also a direct connection with ratio limits for countable
Markov kernels. Restrict the odd-cut transfer \(T\) of
Theorem~\ref{thm:idle-shifts} to the states communicating with \(e\);
every return from \(e\) stays in this class. Then \(T/16\) is
substochastic with diagonal at least \(2/16\). Kesten's lemma gives
the ratio limit for the ordinary meander numbers immediately
\cite[Lemma~4]{K95}; see also \cite[Lemma~1.2]{MSM}.
Deleting only the two specified loops preserves the communicating
class. The returns of lengths two and three in the proof show that
\(Q/16\) is uniformly aperiodic as well, giving the endpoint
ratio by the same lemma. The proof above makes the smoothing
argument explicit without using a general Markov-chain theorem.
